\begin{solutionbox}
\textbf{Given.} $\ket{\psi}=c_1\ket{\phi_1}+c_2\ket{\phi_2}+c_3\ket{\phi_3}$ with $c_1=\frac{\sqrt3}{3}$, $c_2=\frac23$, $c_3=\frac{\sqrt2}{3}$, and $\braket{\phi_i|\phi_j}=\delta_{ij}$.

\textbf{Principle.} Because the basis is orthonormal, $\braket{\psi|\psi}=\sum_i|c_i|^2$ and $c_i=\braket{\phi_i|\psi}$. By the Born rule, a normalised state is found in $\ket{\phi_i}$ with probability $P_i=|\braket{\phi_i|\psi}|^2=|c_i|^2$.

\textbf{(a)} Normalisation:
\[
\braket{\psi|\psi}=\Big(\frac{\sqrt3}{3}\Big)^2+\Big(\frac23\Big)^2+\Big(\frac{\sqrt2}{3}\Big)^2=\frac39+\frac49+\frac29=\frac99=1,
\]
so $\ket{\psi}$ is normalised and no division by the norm is needed. The probabilities are
\[
P_1=\frac39=\frac13,\qquad P_2=\frac49,\qquad P_3=\frac29,
\]
and their total is $\frac39+\frac49+\frac29=1$, as it must be for a complete set of mutually exclusive outcomes.

\textbf{(b)} For $\mathcal N=810$ identically prepared systems, the expected number found in $\ket{\phi_i}$ is $\mathcal N P_i$:
\[
810\times\tfrac13=270,\qquad 810\times\tfrac49=360,\qquad 810\times\tfrac29=180,
\]
and $270+360+180=810$, so every system is accounted for. These are expected (average) counts; an individual run of 810 measurements fluctuates around them by roughly $\sqrt{\mathcal N P_i(1-P_i)}\approx13$--$14$ systems, which a Monte Carlo simulation of 20\,000 such ensembles confirmed.
\end{solutionbox}

\begin{answerbox}
\begin{gather*}
\braket{\psi|\psi}=1;\qquad P(\phi_1)=\tfrac13,\quad P(\phi_2)=\tfrac49,\quad P(\phi_3)=\tfrac29\quad(\text{sum}=1);\\
N(\phi_1)=270,\quad N(\phi_2)=360,\quad N(\phi_3)=180\quad\text{out of }810.
\end{gather*}
Therefore, the state is normalised, the system is found in $\ket{\phi_1},\ket{\phi_2},\ket{\phi_3}$ with probabilities $1/3$, $4/9$ and $2/9$, and in an ensemble of 810 systems we expect $270$, $360$ and $180$ to be found in those states respectively.
\end{answerbox}

\begin{insightbox}
Quantum mechanics predicts the statistics of an ensemble, not the outcome of a single measurement: each of the 810 systems is in exactly the same state, yet they give different results. The ensemble counts in (b) are the operational meaning of the probabilities in (a), and the choice of 810 (divisible by 9) makes the expected counts whole numbers.
\end{insightbox}
